\section{Entropy and induced-graph locality}\label{sec:prelim}

We record the finite-dimensional conventions and elementary uniform bounds
used throughout the proof. Unless specified otherwise, entropies refer to the
original ground vector $\Omega$. A changed state is shown by a subscript.
Juxtaposition of disjoint systems denotes their tensor product or their union.

\subsection{Entropy identities and continuity}

For a density matrix $\rho$ on systems $D,E,F$, define
\[
 S(D|E)_\rho=S(DE)_\rho-S(E)_\rho,
 \quad I(D:E)_\rho=S(D)_\rho+S(E)_\rho-S(DE)_\rho,
\]
and
\[
 I(D:F|E)_\rho=S(DE)_\rho+S(EF)_\rho-S(E)_\rho-S(DEF)_\rho.
\]
We use strong subadditivity, $I(D:F|E)_\rho\ge0$, for normalized
finite-dimensional density operators~\cite{LiebRuskai1973}.
It implies that conditioning decreases conditional entropy and that discarding
systems cannot increase mutual information. Subadditivity and purification
give $|S(D|E)|\le\log\dim D$. For a pure state the entropies of complementary
systems coincide. These facts also hold with empty regions, whose
Hilbert spaces are the one-dimensional empty tensor product.

Write $h(u)=-u\log u-(1-u)\log(1-u)$ for binary entropy.
The trace distance of two states is half the trace norm of their difference.
The continuity estimate below uses the dimension-independent conditional
entropy method of Alicki and Fannes~\cite{AlickiFannes2004}, in Winter's
sharpened form~\cite{Winter2016}.

\begin{lemma}[Continuity with a fixed controlled system]\label{lem:continuity}
Let $\rho,\sigma$ be density operators on finite-dimensional $D\otimes E$,
and put $\delta=\tfrac12\norm{\rho-\sigma}_1\le1$ and $d_D=\dim D$.
Then
\begin{equation}\label{eq:conditional-continuity}
 |S(D|E)_\rho-S(D|E)_\sigma|
 \le2\delta\log d_D+(1+\delta)
                 h\left(\frac{\delta}{1+\delta}\right)
 \le C\sqrt\delta\log(e d_D).
\end{equation}
The same estimates hold for $S(D)$ by making $E$ trivial. Summing the
bounds for $S(D)$ and $S(D|E)$ therefore bounds the continuity cost for
$I(D:E)$ in terms of $d_D$, independently of $\dim E$.
\end{lemma}
\begin{proof}
We include the mixture argument for the bound in
\cite[Lemma~2]{Winter2016}. If $\tau=\sum_j p_j\tau_j$, adding a classical
register that records $j$ shows
\[
 0\le S(D|E)_\tau-\sum_jp_jS(D|E)_{\tau_j}\le H(p).
\]
For the lower bound, condition on that register and use strong
subadditivity. For the upper bound, use the ordinary mixing upper bound on
$S(DE)$ and concavity of $S(E)$. The ordinary bound follows by recording a
duplicate of the classical register: its conditional entropy given $DE$
is nonnegative after this further conditioning.

If $\delta>0$, write the Jordan decomposition
$\rho-\sigma=\delta(\tau_+-\tau_-)$, where $\tau_\pm$ are states.
Then
\[
 \frac{\rho+\delta\tau_-}{1+\delta}
 =\frac{\sigma+\delta\tau_+}{1+\delta}.
\]
Apply the mixture bounds to both expressions and use
$|S(D|E)_{\tau_+}-S(D|E)_{\tau_-}|\le2\log d_D$.
This proves the first inequality. The second follows from the elementary
bound $h(u)\le C\sqrt u$. The case $\delta=0$ is immediate.
Finally, $I(D:E)=S(D)-S(D|E)$, so the two displayed entropy bounds apply.
\end{proof}

The trace distance between pure states represented by unit vectors $\xi,\eta$
is at most $\norm{\xi-\eta}$. We will use Lemma~\ref{lem:continuity} on
regions with a controlled number of sites, or with the controlled region as
the first argument of a conditional entropy. In this way estimates never
acquire the dimension of an uncontrolled complement.

\subsection{Local counting and degenerate cases}

For a site $x$ let $B_\Lambda(x,r)=\{y:d_\Lambda(x,y)\le r\}$.
A path of length $r$ has ambient $\ell^1$ displacement at most $r$; hence,
for an integer $r\ge0$,
\begin{equation}\label{eq:ball-count}
 |B_\Lambda(x,r)|\le v_r:=1+2r(r+1).
\end{equation}
For a physical set $D$ put
$N_r^\Lambda(D)=\{x:d_\Lambda(x,D)\le r\}$.
We reserve graph neighborhoods for support and propagation estimates.
An ambient dilation of a finite set $T\subset\Z^2$ is instead
$T_j=T+([-j,j]^2\cap\Z^2)$ for integer $j\ge0$.

Assign each nonzero interaction in \eqref{eq:hamiltonian} an anchor in its
support. The number of terms containing any prescribed site is at most
\[
 \mu_R:=2^{v_R-1},
\]
because each such support is a subset of its radius-$R$ graph ball containing
that site. In particular, the number of anchored terms at a site is at most
$\mu_R$, and the sum of norms of terms containing that site is at most
$J\mu_R$. Terms with scalar identity factors still obey these bounds with
their assigned supports.

Call an assigned support \emph{split} by a set $D$ when it meets both $D$ and
$\Lambda\setminus D$. For a split support, a graph path of length at most $R$
from its anchor to a point of the opposite side crosses $\partial_\Lambda D$.
The anchor is therefore within graph distance $R$ of a crossing-edge
endpoint. Consequently,
\begin{equation}\label{eq:crossing-count}
 \#\{\text{terms split by }D\}
 \le 2v_R\mu_R|\partial_\Lambda D|.
\end{equation}
Across any bipartition, a term can be decomposed by matrix units on its finite
support into a bounded number of product operators, with bounded factors.
For the original interactions the bounds depend only on $q,R,J$.

If $q=1$ or $\Lambda=\varnothing$, all entropies vanish. When $R=0$, the
Hamiltonian is a sum of on-site operators. Uniqueness of its ground state
then makes that state a tensor product, so the theorem also follows.
We may henceforth assume $q\ge2$, $R\ge1$, and a nonempty domain.

\begin{lemma}[A cut with no crossing edge]\label{lem:zero-boundary}
Under the hypotheses of Theorem~\ref{thm:area}, if
$\partial_\Lambda A=\varnothing$, then $S_\Omega(A)=0$.
\end{lemma}
\begin{proof}
Such a set $A$ is a union of graph components. Every support of finite graph
diameter lies in one component, so $H$ is the tensor sum of component
Hamiltonians. Its ground eigenspace is the tensor product of their ground
eigenspaces. If the former is one-dimensional, so is each latter space.
The ground vector factors over the components and therefore over the cut.
\end{proof}

All constants denoted by $C,c>0$ below may change at different occurrences.
Unless a lemma specifies additional fixed data, they depend only on
$q,R,J,\Delta$. They never depend on the domain, the cut, or the particular
Hamiltonian. Additional numerical exponents and sufficiently large thresholds
will be fixed in the order stated in the proof.
